% Example presentation using cherwellmaths class
% Build with: build.bat example
\documentclass{cherwellmaths}

\title{Introduction to Quadratics}
\subtitle{Year 10 Mathematics}
\author{A. Teacher}

\begin{document}

\section{What is a Quadratic?}

\begin{frame}{General Form}
  \begin{defn}{Quadratic Function}{quad}
    A \emph{quadratic function} has the form
    \[
      f(x) = ax^2 + bx + c, \quad a \neq 0.
    \]
  \end{defn}

  \pause
  Key features:
  \begin{itemize}
    \item Parabola shape (opens up if $a > 0$, down if $a < 0$)
    \item Vertex (turning point)
    \item Axis of symmetry at $x = -\dfrac{b}{2a}$
  \end{itemize}
\end{frame}

\section{Solving Quadratics}

\begin{frame}{The Quadratic Formula}
  \begin{thm*}{Quadratic Formula}{qf}
    If $ax^2 + bx + c = 0$ with $a \neq 0$, then
    \[
      x = \frac{-b \pm \sqrt{b^2 - 4ac}}{2a}.
    \]
  \end{thm*}

  \pause
  \begin{exmp}{Solve $x^2 - 5x + 6 = 0$}{ex1}
    Here $a=1$, $b=-5$, $c=6$.
    \[
      x = \frac{5 \pm \sqrt{25 - 24}}{2} = \frac{5 \pm 1}{2}
    \]
    So $x = 3$ or $x = 2$.
  \end{exmp}
\end{frame}

\section{Graphing}

\begin{frame}{Graph of $y = x^2 - 5x + 6$}
  \centering
  \begin{tikzpicture}
    \begin{axis}[
      axis lines=middle,
      xlabel=$x$, ylabel=$y$,
      xmin=-1, xmax=6, ymin=-2, ymax=7,
      width=\slidesquare,
      height=0.8\slidesquare,
      grid=major,
    ]
      \addplot[domain=-0.5:5.5, samples=80, thick, color=cherwellaccent]{x^2 - 5*x + 6};
      \addplot[only marks, mark=*, color=cherwellblack] coordinates {(2,0)(3,0)(2.5,-0.25)};
    \end{axis}
  \end{tikzpicture}
\end{frame}

\section{Side-by-Side Demo}

\begin{frame}{Graph with Key Points}
  \sidebyside{%
    \begin{tikzpicture}
      \begin{axis}[
        axis lines=middle,
        xlabel=$x$, ylabel=$y$,
        xmin=-1, xmax=6, ymin=-2, ymax=7,
        width=\slidesquare,
        height=\slidesquare,
        grid=major,
      ]
        \addplot[domain=-0.5:5.5, samples=80, thick, color=cherwellaccent]{x^2 - 5*x + 6};
      \end{axis}
    \end{tikzpicture}
  }{%
    Key features of $y = x^2 - 5x + 6$:
    \begin{itemize}
      \item Roots at $x = 2$ and $x = 3$
      \item Vertex at $(2.5,\, -0.25)$
      \item Opens upward ($a > 0$)
      \item $y$-intercept at $(0, 6)$
    \end{itemize}
  }
\end{frame}

\begin{frame}{Key Points with Graph}
  \sidebyside{%
    To factorise $x^2 - 5x + 6$:
    \begin{itemize}
      \item Find two numbers that multiply to $6$
      \item \ldots and add to $-5$
      \item $-2 \times -3 = 6$ \checkmark
      \item $-2 + (-3) = -5$ \checkmark
      \item So $y = (x-2)(x-3)$
    \end{itemize}
  }{%
    \begin{tikzpicture}
      \begin{axis}[
        axis lines=middle,
        xlabel=$x$, ylabel=$y$,
        xmin=-1, xmax=6, ymin=-2, ymax=7,
        width=\slidesquare,
        height=\slidesquare,
        grid=major,
      ]
        \addplot[domain=-0.5:5.5, samples=80, thick, color=cherwellaccent]{x^2 - 5*x + 6};
      \end{axis}
    \end{tikzpicture}
  }
\end{frame}

\section{Incremental Reveal Demo}

\begin{frame}{Building Up a Solution}
  Solve $x^2 - 4x + 3 = 0$ by factorising:

  \pause
  \textbf{Step 1:} Find two numbers that multiply to $3$ and add to $-4$.

  \pause
  \textbf{Step 2:} The numbers are $-1$ and $-3$.

  \pause
  \textbf{Step 3:} Write the factorised form:
  \[
    (x - 1)(x - 3) = 0
  \]

  \pause
  \textbf{Step 4:} So $x = 1$ or $x = 3$.
\end{frame}

\begin{frame}{Revealing Bullet Points One by One}
  Properties of quadratic graphs:
  \begin{itemize}[<+->]
    \item Always a parabola
    \item Exactly one turning point (vertex)
    \item Axis of symmetry passes through the vertex
    \item At most two $x$-intercepts
    \item Exactly one $y$-intercept
  \end{itemize}
\end{frame}

\end{document}
